\section{Cross-organization admission}
\label{sec:cross}

A call from another organization is the same admission operation. The peer
kernel's receipt and co-signature enter at step 4 as evidence and add no
authority: the receiver's issuer set, pinned peers, agreements, and clocks are
what they were before the request arrived. Figure~\ref{fig:mechanism} is the
whole of what a statement can name.

\subsection{Receiver-owned resolution}

Two organizations pin each other out of band, each recording the other's kernel
identity, the passport key that kernel co-signs with, and a rotation deadline.
The agreement is an activated record naming exactly two kernels, the action
classes in scope, a validity window, and a revocation epoch. Each participant
publishes a signed manifest, the action classes it accepts and at what
intensity on the ascending ladder \code{observation}, \code{guarded},
\code{receipt_backed}, \code{partition_contingency},
\code{maintenance}, and each receiver activates the agreement and
both manifests locally and intersects the two per action class itself: the
higher-intensity mode wins, the stricter co-signing requirement wins, a
destructive class below \code{receipt_backed} rejects, and
mismatched consistency models reject. Consistency is a separate axis, not a
mode. Continuations resolve by identifier against the receiver's own copy. A pin the receiver has not activated is a
denial, not a prompt.

The request may carry identifiers. It may not carry trust. Its agreement context
is names and digests, each a lookup key into a store the receiver owns: the
agreement and the intersection with the digest the receiver should find for
each, the action class, and identifiers for the continuation, the lineage
bundle that chains this call to the receipts behind it, and the co-signed
evidence. Eleven further key names are refused wherever they appear. Seven name
a trust root, a bundle, an agreement record, a manifest, a signing key, or a
peer directory, refused with
\code{request_smuggled_trust_root}: \code{trustRoot}, \code{trustRoots},
\code{trustBundle}, \code{treatyScope}, \code{ladderManifest},
\code{signingKey}, \code{peerDirectory}. Four name dynamic trust input or peer
discovery, refused with \code{request_smuggled_dynamic_trust}:
\code{dynamicTrust}, \code{dynamicTrustBundle}, \code{runtimeTrustInput},
\code{peerDiscovery}. The check is on names, runs before anything is resolved,
and is the operational form of the rule.

\subsection{The co-signed statement}

The statement is an in-toto Statement v1 \cite{intoto} in a DSSE envelope
\cite{dsse} whose predicate is snake\_case and rejects unknown fields. Its
single subject names the invocation identifier, which is the receipt's
\code{id}, and carries as its digest the SHA-256 of the canonical receipt body
of Section~3. The envelope for the refund step, then its payload, abridged from
the committed evidence package for workflow \code{wf-chio-refund-001}:
\begin{lstlisting}
{ "payloadType": "application/vnd.in-toto+json",
  "payload": "eyJfdHlwZSI6Imh0dHBzOi8vaW4tdG90by5pby9T...",
  "signatures": [
    { "keyid": "fdf72a088f18f7399e8c52bce448441501f759a595a86980e5d9a422a01e5d55",
      "sig": "3QZ09JP4YLFmvgzwy5PWqYTWFCcQQkUo/DRC7YyFB8pD+bA6iBGwyRTtI+lzzI7A..." },
    { "keyid": "e37abcd79efd08db94518a903952ffe9f8af7e8b4684b63591631cef8e8871c8",
      "sig": "R0Ii8tcBzYcEPAPRQqNmsbCUowFH0hzspjr5Mm98mFu33+XcT0fDxdWf48btYzmd..." } ] }
\end{lstlisting}

\begin{lstlisting}
{ "_type": "https://in-toto.io/Statement/v1",
  "predicateType": "chio.bilateral-cosign-invocation.v1",
  "subject": [ { "name": "chio-receipt:d342a59d...6fba2e3f",
                 "digest": { "sha256": "31118bb6...4781b3f9" } } ],
  "predicate": {
    "invocation_id": "d342a59d...6fba2e3f", "tool_name": "stage_refund",
    "tool_server_a": { "kernel_id": "did:chio:buyer-kernel", "alg": "ed25519",
      "passport_key_fingerprint": "fdf72a08...a01e5d55" },
    "tool_server_b": { "kernel_id": "did:chio:vendor-c", "alg": "ed25519",
      "passport_key_fingerprint": "e37abcd7...8e8871c8" },
    "co_sign": "bilateral_required", "consistency_model": "totally-ordered",
    "consistency_anchor": "chio:consistency:wf-chio-refund-001:2",
    "cross_org_visibility": "federated", "timestamp_unix_ms": 1766000001000,
    "tool_args_hash": { "alg": "sha256", "value": "47e6e096...d5a6464b" },
    "capability_lease_ref": { "lease_id": "lease-vendor-c-refund",
      "issuer": "did:chio:buyer-kernel", "expires_at_unix_ms": 1766000030000,
      "scope_digest": { "alg": "sha256", "value": "262768fe...b9f835aa" } },
    "governance_receipt_ref": { "receipt_id": "gov-refund-stage-authorization",
      "kernel_id": "did:chio:buyer-governance",
      "digest": { "alg": "sha256", "value": "21143a93...f161e3c0" } },
    "policy_evaluation_summary": { "joint_disposition": "allow",
      "server_a_verdict": { "verdict": "allow", "policy_id": "buyer-policy:stage_refund", ... },
      "server_b_verdict": { "verdict": "allow", ... } },
    "treaty_binding_ref": {
      "treaty_id": "treaty:runtime-loopback:2",
      "action_class_id": "workflow.cross_kernel.stage_refund",
      "consistency_model": "totally-ordered",
      "treaty_scope_sha256":        "0ff230e9...d1f2b21c",
      "ladder_intersection_sha256": "93dcc7cf...af79d3d1",
      "admission_report_sha256":    "91d74509...1f646522",
      "continuation_sha256":        "8ef2f348...366b93f4",
      "lineage_bundle_sha256":      "28d816a3...d3b504ba",
      "request_sha256":             "47e6e096...d5a6464b",
      "outcome_sha256":             "c9abf68a...d9803987",
      "local_receipt_sha256":       "9cbe4438...ad017121",
      "remote_receipt_sha256":      "63f1f35a...8c93e3e8",
      "lease_refs":        [ "lease-vendor-c-refund" ],
      "governance_refs":   [ "gov-refund-stage-authorization" ],
      "signer_kernel_ids": [ "did:chio:buyer-kernel", "did:chio:vendor-c" ] } } }
\end{lstlisting}

Both signatures cover the DSSE pre-authentication encoding of the payload, which
here begins \code{DSSEv1 28 application/vnd.in-toto+json 2972 \{"_type":"https:}
and runs to 3016 bytes; over those bytes both verify, one under each
kernel's key, and each \code{keyid} is the SHA-256 of its raw public key.
Every identity above was recomputed from the committed bytes:
\code{outcome_sha256} is the receipt's \code{content_hash}, and
\code{request_sha256} is both \code{tool_args_hash.value} and the receipt's
\code{action.parameter_hash}.

\subsection{The verifier, in the order it runs}

Table~\ref{tab:verifier} is the shipped order. Two components decide, in two
disjoint code families. The conforming verifier returns the
specification's dotted codes, a closed set of \PSRejectionCodeCount{} strings of
which it emits \PSVerifierCodeCount{} and the envelope layer the rest. The
kernel's pre-dispatch hook, which resolves the same evidence inline during
admission, answers from the kernel's runtime failure-code set of
\PSRuntimeCodeCount{} strings, \PSRuntimeTreatyCodeCount{} of them
agreement-scoped, and maps every envelope-layer failure to
\code{chio_treaty_unverified_required_evidence}. Table~\ref{tab:subst} and the
denial experiment of Section~\ref{sec:eval} report the hook, because the hook
decides a live call.
\begin{table}[!tb]
\centering
\scriptsize
\caption{The conforming verifier, \PSVerifierStepCount{} steps. The receiver's
own state decides everything outside steps 8 to 14, which are the envelope
layer, called with the two keys the receiver pinned.}
\label{tab:verifier}
\begin{tabular}{@{}rp{7.4cm}p{5.5cm}@{}}
\toprule
\# & Check & Code on failure \\
\midrule
1--5 & Media type, at least one signature, payload decodes and parses; predicate type \code{chio.\allowbreak{}bilateral-cosign-\allowbreak{}invocation.v1}; Statement v1; exactly one subject & \code{dsse.malformed}, \code{statement.malformed}, \code{predicate.type_unrecognised}, \code{statement.schema_invalid} \\
6--7 & Both kernel identifiers are in the pinned peer set and each pin's fingerprint equals the declared one & \code{peer.unpinned_or_keyid_mismatch} \\
8--13 & Exactly two signatures; re-canonicalizing reproduces the payload bytes; predicate schema and strict rules; both verdicts present, equal, in $\{\mathrm{allow},\mathrm{deny}\}$; subject name is \code{chio-receipt:} plus the invocation identifier; the two pinned keys and their identifiers distinct & \code{statement.malformed}, \code{predicate.schema_invalid}, \code{policy.verdict_disagreement}, \code{subject.digest_mismatch}, \code{signer.independence_required} at the envelope layer, surfacing as \code{dsse.malformed} \\
14 & \emph{Only now:} compute the pre-authentication encoding, select each signature by a pinned key's identifier, and verify both & \code{signature.server_a_invalid}, \code{signature.server_b_invalid} \\
15--16 & Each pinned peer's manifest reference is present and fresh, and the peer is active at the pinned epoch height & \code{ladder.manifest_missing}, \code{ladder.manifest_stale}, \code{peer.revoked_at_epoch} \\
17--19 & Resolve the receipt by invocation identifier in the receiver's own store; its signature verifies, its tool name and argument hash agree, its kernel key is the pinned second-party key, and subject name and digest match the resolved body & \code{subject.digest_mismatch}, \code{peer.unpinned_or_keyid_mismatch} \\
20--22 & Both verdicts resolve and agree; the lease resolves with matching issuer, expiry and scope digest and is unexpired; the tool name is in the action-class table; a receipt-backed class resolves its governance record & \code{policy.verdict_disagreement}, \code{capability.lease_expired_or_unknown}, \code{governance.unknown_action_class}, \code{governance.receipt_required_missing} \\
23--26 & Binding cross-checks (Table~\ref{tab:subst}), consistency model and anchor, co-signing mode; return the verified material & \code{predicate.schema_invalid} \\
\bottomrule
\end{tabular}
\end{table}

Signatures are verified last and the pin lookup first, the opposite of the usual
advice: the verifier cannot verify until it knows which keys to verify under,
and it takes those from its own pin set, never from the envelope, so a signature
whose key identifier matches no pinned key is a missing signature. No gate before
step 14 consults a secret, so the ordering costs nothing in soundness. The typed error exposes a code and keeps the rendered detail for the
local log, since that message names presented and expected digests and would
make the verifier an oracle.

\begin{figure}[!tb]
\centering
\begin{tikzpicture}[
  font=\small,
  env/.style={draw,rounded corners=2pt,align=center,inner sep=3pt,fill=white},
  rec/.style={draw,rounded corners=2pt,align=left,inner sep=2pt,fill=black!3},
  ptr/.style={-{Latex[length=1.6mm]},semithick},
]
\node[env,align=left,text width=42mm] (stmt) at (0,0)
  {\textbf{co-signed statement}\\[1pt]
   \scriptsize subject: one name, one SHA-256\\
   \scriptsize binding reference: \PSBindingFieldCount{} fields\\
   \scriptsize signatures: exactly 2, same bytes};
\node[rec] (rcpt) at (58mm,12mm) {\scriptsize a receipt already in the store};
\node[rec] (scope) at (58mm,6.5mm) {\scriptsize the activated agreement and its intersection};
\node[rec] (cont)  at (58mm,1mm)  {\scriptsize an unspent continuation};
\node[rec] (lease) at (58mm,-4.5mm) {\scriptsize a live lease, a governance record};
\node[rec] (keys)  at (58mm,-10mm) {\scriptsize two separately pinned peer keys};
\node[rec] (epoch) at (58mm,-15.5mm) {\scriptsize the revocation clock, fresh};
\foreach \t in {rcpt,scope,cont,lease,keys,epoch}{
  \draw[ptr] (stmt.east) .. controls +(11mm,0) and +(-9mm,0) .. (\t.west);
}
\node[font=\scriptsize\itshape,anchor=south] at (29mm,15mm)
  {every arrow terminates inside the receiver};
\end{tikzpicture}
\caption{The mechanism. The statement carries no authority of its own: its
subject, the binding fields the receiver compares, and each of its two keys
resolve only against records the receiving kernel placed in its own stores
before the request arrived.}
\label{fig:mechanism}
\end{figure}


\subsection{Consume, bind, co-sign}

Replay is not defeated by a signature, since a replayed statement is validly
signed. It is defeated by a single-use continuation: a token minted by the
origin binding the parent receipt, the parent session anchor, the capability,
the action class, the target tool, a nonce, the two kernel identities, and a
validity window, whose digest is one of the \PSBindingFieldCount{} bound fields.
The receiver's state machine over it is two operations on one table keyed by the
continuation identifier: \emph{consume} inserts, ignoring a conflict, and reads
the affected row count; zero rows means the continuation was already spent and
the call is denied with \code{chio_treaty_continuation_replay}. \emph{Release}
deletes the row for that identifier and that admission, and runs only when a
later pre-dispatch step denies, so a rejected call does not burn a continuation
the origin will retry. Once dispatch commits the row stays: at-most-once, on
purpose.

The receiver then builds its own admission report and binds the accepted
material into the receipt it is about to sign. One field is treated differently
from the other fourteen: the receiver checks only that
\code{admission_report_sha256} is well-formed hex, then recomputes the digest of its \emph{own} report and overwrites the field before
signing, because a peer-supplied hash must never sit inside a locally signed
receipt as though the receiver had computed it. Verified material is installed
for that one request as a validated value, never as configuration. The local
receipt is persisted; only then does the kernel ask the peer to co-sign a
statement whose subject is the new receipt's digest, so a refusal or a transport
failure is recorded against a receipt that already exists rather than erasing
it.
